\documentclass{article}
\usepackage{graphicx} % Required for inserting images
\usepackage{xspace}
\usepackage{amsmath, amssymb}

\newcommand{\bred}{$\beta$-reduction\xspace}
\newcommand{\lcalc}{$\lambda$-calculus\xspace}
\newcommand{\lterm}{$\lambda$-term\xspace}
\newcommand{\lterms}{$\lambda$-terms\xspace}


\title{Notes on Theories}
\author{Clément Michaud}
\date{January 2025}

\begin{document}

\maketitle

\section{Untyped \lcalc}

\subsection{Simple construction}

It is a way to abstract function on anything. Therefore, one does not need any notion of algebra or whatever. What remains is just:

\begin{itemize}
    \item variables
    \item creating an expression representing an abstract function with variables (Abstraction)
    \item apply the argument to create a new expression (it’s not the execution of the function) (Application)
    \item \bred is the execution of the function.

\end{itemize}

More formally, we assume the existence of an infinite set of variables $V=\{x, y, z, \dots\}$. Expressions in the \lcalc are called \lterms. The set of all \lterms is $\Lambda$.\\

\textbf{Definition} (The set $\Lambda$ of all $\lambda$-terms)
\begin{enumerate}
    \item (Variable) If $u \in V$, then $u \in \Lambda$.
    \item (Application) If $M, N \in \Lambda$, then $(MN) \in \Lambda$.
    \item (Abstraction) If $u \in V$ and $M \in \Lambda$, then $(\lambda u . M) \in \Lambda$.
\end{enumerate}

\subsection{Theorems}

\noindent \textbf{Proposition:} If $L \in \text{Sub}(M)$, then $\text{Sub}(L) \subseteq \text{Sub}(M)$.

\bigskip

\noindent \textbf{Proof:} We prove the proposition by structural induction on the $\lambda$-term $M$.

\medskip

\noindent \textbf{Base Case:} Let $M$ be a variable $x \in V$. By definition, $\text{Sub}(x) = \{x\}$. If $L \in \text{Sub}(x)$, the only possibility is $L = x$, and thus $\text{Sub}(L) = \text{Sub}(x) = \{x\} \subseteq \text{Sub}(x)$. Therefore, the base case holds.

\medskip

\noindent \textbf{Inductive Case 1:} Suppose $M = (PQ)$ for some $\lambda$-terms $P$ and $Q$. By definition:
\[
\text{Sub}((PQ)) = \text{Sub}(P) \cup \text{Sub}(Q) \cup \{(PQ)\}
\]
Assume that the proposition holds for all subterms of $P$ and $Q$.

We need to prove that if $L \in \text{Sub}((PQ))$, then $\text{Sub}(L) \subseteq \text{Sub}((PQ))$.

There are three possibilities for $L$:
\begin{enumerate}
    \item If $L \in \text{Sub}(P)$, then by the induction hypothesis, $\text{Sub}(L) \subseteq \text{Sub}(P)$. Since $\text{Sub}(P) \subseteq \text{Sub}((PQ))$, we have $\text{Sub}(L) \subseteq \text{Sub}((PQ))$.
    \item If $L \in \text{Sub}(Q)$, then by the induction hypothesis, $\text{Sub}(L) \subseteq \text{Sub}(Q)$. Since $\text{Sub}(Q) \subseteq \text{Sub}((PQ))$, we have $\text{Sub}(L) \subseteq \text{Sub}((PQ))$.
    \item If $L = (PQ)$, then $\text{Sub}(L) = \text{Sub}((PQ))$, which trivially implies $\text{Sub}(L) \subseteq \text{Sub}((PQ))$.
\end{enumerate}

Thus, the inductive case for application holds.

\medskip

\noindent \textbf{Inductive Case 2:} Suppose $M = (\lambda x . P)$ for some $\lambda$-term $P$. By definition:
\[
\text{Sub}((\lambda x . P)) = \text{Sub}(P) \cup \{(\lambda x . P)\}
\]
Assume that the proposition holds for all subterms of $P$.

We need to prove that if $L \in \text{Sub}((\lambda x . P))$, then $\text{Sub}(L) \subseteq \text{Sub}((\lambda x . P))$.

There are two possibilities for $L$:
\begin{enumerate}
    \item If $L \in \text{Sub}(P)$, then by the induction hypothesis, $\text{Sub}(L) \subseteq \text{Sub}(P)$. Since $\text{Sub}(P) \subseteq \text{Sub}((\lambda x . P))$, we have $\text{Sub}(L) \subseteq \text{Sub}((\lambda x . P))$.
    \item If $L = (\lambda x . P)$, then $\text{Sub}(L) = \text{Sub}((\lambda x . P))$, which trivially implies $\text{Sub}(L) \subseteq \text{Sub}((\lambda x . P))$.
\end{enumerate}

Thus, the inductive case for abstraction holds.

\medskip

\noindent \textbf{Conclusion:} By structural induction on $M$, we have shown that if $L \in \text{Sub}(M)$, then $\text{Sub}(L) \subseteq \text{Sub}(M)$.

\hfill $\Box$

\subsection{Grammar}

$\Lambda = V \mid (\Lambda\Lambda) \mid (\lambda V.\Lambda) $

\end{document}
